\section{What the executable core proves}
\label{sec:v2proof}
The \texttt{policy-core} crate is a dependency of the broker and sink. Its
Verus-annotated source is also the source compiled by ordinary Cargo after ghost
erasure. There is no separate proof-only policy implementation. The verifier
release archive, archive hash, Rust toolchain, dependency versions, and source
hash are recorded. Verification runs both directly on the core and through
Cargo's Verus integration with the locked dependency graph \cite{verus}.

\subsection{Decision specifications}
The core's byte comparison proves equivalence to equality of the complete
finite byte sequences. This connects positive and negative permission results:
identical sequences are accepted and any differing byte or length is rejected.
Parsing and normalization must still supply the intended sequences. The proof
does not assert that arbitrary JSON objects, natural-language descriptions, or
two external tool names denote the same effect.

The grant predicate proves the conjunction over all required entries in a
finite vector. A second executable predicate additionally checks liveness and
equality of current and requested epochs for each required issuer. Length
mismatches fail closed. Nonrequired issuers do not create a spurious veto.
These properties quantify over vector length rather than enumerating two
actors or a fixed list of interleavings. The application supplies the required
vector and binds it to the normalized effect; that mapping is trusted code.

The admission predicate is specified by an equivalence, not only the implication
``if accepted, then safe.'' It accepts exactly when the budget, expiry,
liveness, epoch, revision, payload, and required-grant conditions in its inputs
hold. Some application checks happen before admission or at the sink: issuer
epochs are checked by the current-grants predicate, and the sink checks actual
document revision. The broker passes neutral placeholders for revision in its
earlier admission call because the authoritative object is in the sink's store.
The proof therefore must not be read as an atomic cross-store revision check.

The reservation function returns either no reservation or the exact new charge.
Its postconditions prove positivity, capacity, and absence of arithmetic
overflow. The release function never increases charged authority and leaves it
unchanged without authoritative no-effect evidence. The authoritative nature of
that evidence is a caller obligation. A Boolean in verified code cannot
authenticate the message that supplied it.

The fence predicate accepts exactly when every required request epoch is at
least its corresponding sink floor. Floors are advanced monotonically by the
storage adapter. A separate lemma shows that an old epoch for any one required
issuer makes the universal acceptance predicate false. Revision equality is
also an executable verified predicate. The sink uses these predicates within
the same transaction as effect application and terminal-status persistence.

\subsection{Conservation across finite prefixes}
The arithmetic proof is independent of which actor proposes the next request.
Represent charged authority after each step by $C_t$, with $C_0\in[0,B]$.
A reservation has increment $d_t>0$ and requires $d_t\le B-C_t$.
A release has $d_t\le0$ and requires $-d_t\le C_t$; recovery,
renewal, and terminal committed settlement have $d_t=0$.
For $C_{t+1}=C_t+d_t$, induction gives
\begin{equation}
                 0\le C_t\le B
\label{eq:v2conservation}
\end{equation}
for every finite prefix satisfying those transition premises. The mechanized
prefix lemma performs this induction over arbitrary finite sequences. The
executable reservation and release specifications discharge the arithmetic
premises when called with valid state. There is no multiplication by actor
population and no new budget on restart or renewal.

This is a conditional composition argument. The database adapter must serialize
updates, load the correct durable $C_t$, persist the returned $C_{t+1}$, and avoid
unmediated effects. SQLite immediate transactions provide the implementation's
serialization mechanism. SQL execution and crash recovery are tested, not
formally refined into the prefix theorem. An omitted caller, stale storage
snapshot, or direct sink credential would break the correspondence even if
every core theorem continued to pass.

For replay, the stored identity locates an immutable permit. The same verified
byte comparison checks the retry. The broker returns its existing status rather
than charging again; the sink returns terminal status rather than applying
again. At-most-once effect application also depends on the sink's unique identity
constraint and atomic transaction. Those database facts are trusted premises,
with duplicate, altered-request, crash, and restart regressions as implementation
evidence. We do not claim a mechanized SQL at-most-once theorem.

\subsection{Positive permissions and specification omissions}
A reject-all implementation cannot satisfy the admission equivalence when the
right-hand side is true. Runtime positive controls demonstrate authorized
transfers, exact releases, and an edit producing the expected contents. Vector
tests cover independent required and nonrequired issuer behavior, while the
proof itself is not bounded by those test lengths. Positive permission says
the interface can admit specified useful work; it does not establish general
task coverage.

A deliberate omission separates permission from harm. Both issuers approve an
exact edit whose function returns zero, while the task requires 42. The verified
policy admits it and the supported sink commits it. An independent content
oracle fails the task. This is expected: the permission specification contains
exact approval, not a theorem about the user's intended program. The negative
result is retained with the successful proof output rather than rewritten as an
authorization failure.

The same distinction limits what a larger proof could buy. Formalizing all
caller transitions would strengthen implementation correspondence but still
would not establish the adequacy of the policy or legitimacy of its issuers.
The trusted base needs independent review, semantic task oracles, and concrete
boundary probes. Verification and testing answer different questions, and the
release reports their evidence separately.
